% Auto-generated from fsmrepairbench.benchmark_utility_calibration
\begin{table}[t]
\caption{Minimum detectable complete-repair difference (MDE) at 80\% power using a simple two-proportion normal approximation ($\alpha=0.05$, two-sided). All detectable cases ($n=495$): 7.9~pp at observed pooled rate 26.9\%; conservative bound 8.9~pp at pooled rate 50\%. Per-operator cells ($n\approx55$): 26.7~pp. Assumes independent equal-sized samples; paired McNemar tests on the same cohort can be more powerful.}
\label{tab:minimum-detectable-effect}
\begin{tabular}{@{}lrrrp{0.42\linewidth}@{}}
\toprule
Scope & $n$ & Pooled rate & MDE (pp) & Assumption \\
\midrule
All detectable (observed pooled) & 495 & 26.9\% & 7.9 & Independent two-proportion normal approximation with equal sample sizes; alpha=0.05 two-sided; power=0.80; conservative pooled rate shown. \\
All detectable (pooled 0.5) & 495 & 50.0\% & 8.9 & Independent two-proportion normal approximation with equal sample sizes; alpha=0.05 two-sided; power=0.80; conservative pooled rate shown. Pooled rate fixed at 0.5 for conservative upper bound. \\
Per-operator cell (mean $n$) & 55 & 50.0\% & 26.7 & Independent two-proportion normal approximation with equal sample sizes; alpha=0.05 two-sided; power=0.80; conservative pooled rate shown. Uses mean detectable cases per operator cell. \\
\bottomrule
\end{tabular}
\end{table}
